% !TEX root = ../main.tex
\section{Trabajo preparatorio y fundamentos algebraicos}

\subsection{Unidades módulo \texorpdfstring{$n$}{n}, función \texorpdfstring{$\varphi$}{phi} y orden}
El conjunto $\mathbb{Z}_n^{*}=\{a\bmod n:\gcd(a,n)=1\}$ es un grupo multiplicativo de orden $\varphi(n)$. El orden de una unidad $a$ es el menor $t>0$ con $a^t\equiv1\pmod n$; por el teorema de Lagrange, divide a $\varphi(n)$. Euler establece $a^{\varphi(n)}\equiv1\pmod n$ cuando $\gcd(a,n)=1$. Si $n=p$ es primo, $\varphi(p)=p-1$: se obtiene el pequeño teorema de Fermat \cite{shoup,cormen}. La contraposición permite \emph{certificar composición} cuando falla la congruencia; el recíproco no es válido.

\subsection{División de prueba, pseudoprimos y Carmichael}
Si $n=ab$ es compuesto y ambos factores excedieran $\sqrt n$, entonces $ab>n$. Por ello basta comprobar los posibles divisores hasta $\lfloor\sqrt n\rfloor$. Para $L=\lceil\log_2(n+1)\rceil$ bits, esta cota es de orden $2^{L/2}$ divisiones; no es polinómica en $L$. Una base $a$ coprima a un compuesto $n$ que cumple $a^{n-1}\equiv1\pmod n$ convierte a $n$ en un \emph{pseudoprimo de Fermat} para $a$. Por ejemplo, $341=11\cdot31$ pasa la base $2$. Un número de Carmichael pasa todas las bases coprimas: $561=3\cdot11\cdot17$, $1105=5\cdot13\cdot17$ y $1729=7\cdot13\cdot19$ son casos de prueba. El criterio de Korselt caracteriza tales números como compuestos libres de cuadrados con $(p-1)\mid(n-1)$ para todo primo $p\mid n$ \cite{crandall}.

\subsection{Exponenciación binaria y Miller--Rabin}
La representación binaria $e=\sum b_i2^i$ permite computar $a^e\bmod n$ mediante cuadrados sucesivos: $O(\log e)$ productos modulares, sin construir $a^e$ completo. La instrumentación cuenta cuadrados \emph{ejecutados}, incluidos los de la exponenciación inicial y la cadena posterior. Para $n>2$ impar se descompone $n-1=2^sd$ con $d$ impar. La ronda comienza con $x=a^d\bmod n$; acepta si $x=1$ o si $x=-1\pmod n$ al inicio o en alguno de los $s-1$ cuadrados posteriores. Si la cadena alcanza $1$ desde $x\not\equiv\pm1$, se encontró una raíz no trivial de $1$, imposible en un cuerpo $\mathbb{Z}_p$: se certifica composición y $\gcd(x\pm1,n)$ puede proporcionar factores. Cuando una base rechaza al compuesto se denomina \emph{testigo}; cuando lo acepta, \emph{mentiroso fuerte}.

\subsection{Probabilidad de error y tipo de algoritmo}
Para un compuesto impar fijo, como máximo un cuarto de las bases aleatorias válidas son mentirosos fuertes. Con $r$ bases uniformes independientes, la probabilidad de aceptar ese compuesto es a lo sumo $4^{-r}$. Para $r=1,5,10,20$ las cotas son $0.25$, $9.765625\cdot10^{-4}$, $9.536743\cdot10^{-7}$ y $9.094947\cdot10^{-13}$. La conclusión ``compuesto'' es definitiva; ``probable primo'' no es una prueba matemática. Miller--Rabin es Monte Carlo de error unilateral: tiempo acotado por rondas y posible falso positivo, nunca falso negativo para un primo. Un algoritmo Las Vegas siempre entrega una respuesta correcta, pero su tiempo de ejecución puede depender del azar. Las semillas fijas reproducen los experimentos; para usos criptográficos se requiere \texttt{secrets.SystemRandom} y entropía no predecible \cite{pythonsecrets}.

\subsection{Modelo de costos y expectativa de candidatos}
Por ronda, Fermat y Miller--Rabin requieren $O(L)$ productos modulares; con $r$ rondas, Miller--Rabin emplea $O(rL)$ productos, frente a $O(2^{L/2})$ divisiones de prueba. Para multiplicación escolar de enteros de $L$ bits, la estimación de costo por ronda es $O(L^3)$ operaciones de bit. Entre impares de $L$ bits, la densidad aproximada de primos es $2/(L\ln2)$: se esperan alrededor de $L\ln2/2$ candidatos impares por primo. El precribado evita exponenciaciones innecesarias, pero no cambia la probabilidad de que un candidato generado al azar sea primo.
